% Propietario conservado: /Users/ruben/Documents/New project/output/REVISION_ARGUMENTAL_APERTURAS_CIERRES_20260915/VIII/spanish_source/gestion/lectura_generada/02_ORDEN_ESPECTRAL_Y_MOMENTOS.tex
\section{Factorisation, horloges arithmétiques et hiérarchie des traces}
\label{viii:lectura:manuscrito-02}
Le produit natif permet de distinguer deux propriétés : l'irréductibilité d'une forme normale et la période d'une action multiplicative sur celle-ci. Ce chapitre construit leurs lectures spectrales, conserve les normalisations de chaque fonctionnelle et prépare la forme complétée développée dans la section~\ref{viii:lectura:manuscrito-03}.

\subsection{De l'irréductible à son horloge}
Le produit et le coproduit déterminent quels objets sont irréductibles. L'horloge arithmétique décrit une autre propriété de ces objets : le retour d'une orbite multiplicative finie. Pour un entier $n\ge2$ premier avec dix, on définit

\[
\tau_n=\min\{k\ge1:10^k\equiv1\pmod n\}.
\]
La longueur $\tau_n$ correspond à l'orbite de l'unité sous la multiplication par dix dans le groupe des unités modulo $n$. La base décimale fixe cette lecture ; une autre base première avec le module détermine une autre horloge. Le caractère irréductible appartient en revanche au produit construit avant le choix de cette représentation périodique.

\HMTresultado{Proposition}{Synchronisation des horloges}{rz-num-02-1}
Si $n=\prod_p p^{a_p}$ et $\gcd(n,10)=1$, alors

\[
\tau_n=\operatorname{mcm}_{p^{a_p}\parallel n}\tau_{p^{a_p}}.
\]
\textbf{Démonstration.} Un exposant $k$ satisfait $10^k\equiv1\pmod n$ si et seulement s'il la satisfait modulo chaque puissance première de la factorisation. Cela équivaut à ce que $k$ soit un multiple de chaque ordre $\tau_{p^{a_p}}$. Le plus petit exposant positif commun est le plus petit commun multiple. \hfill\(\square\)

La factorisation organise donc une synchronisation. Les puissances d’un même nombre premier doivent être maintenues comme telles : leur période ne s’obtient pas en les remplaçant sans preuve par la période du nombre premier de base.

\HMTresultado{Proposition}{Clôture résiduelle nonadique}{rz-num-02-2}
Si $\gcd(n,90)=1$, le bloc entier

\[
M_n=\frac{10^{\tau_n}-1}{n}
\]
est divisible par neuf.

\textbf{Démonstration.} La définition de la période rend $M_n$ entier. Puisque $10\equiv1\pmod9$, le numérateur est divisible par neuf. L'entier $n$ est inversible modulo neuf, de sorte que $nM_n\equiv0\pmod9$ implique $M_n\equiv0\pmod9$. \hfill\(\square\)

Pour un nombre premier $p\notin\{2,3,5\}$, on a en outre $\tau_p\mid p-1$, puisque son orbite appartient au groupe multiplicatif d’ordre $p-1$. La congruence de la proposition~\ref{viii:rz-num-02-2} est également satisfaite pour les composés premiers avec quatre-vingt-dix. Elle signifie une clôture résiduelle ; l’irréductibilité demeure déterminée par les ouvertures du produit.

Les nombres $7,13,77,91,143$ ont une période décimale de six, mais seuls les deux premiers sont irréductibles. La relation structurelle ne consiste pas en ce qu'une période identifierait à elle seule un nombre premier, mais en ce que le produit construit des composantes dont la composition organise les retours.

\subsection{Factorisation comme occupation de modes}
Soit $\mathcal P$ l'ensemble obtenu en évaluant les irréductibles natifs. Chaque occupation de modes est une famille $\mathbf k=(k_p)_{p\in\mathcal P}$ d'entiers non négatifs à support fini. Son évaluation multiplicative est

\[
N(\mathbf k)=\prod_{p\in\mathcal P}p^{k_p}.
\]
La factorisation unique identifie ces occupations aux entiers positifs. L’espace de Hilbert des occupations et celui des formes normales positives possèdent alors une correspondance unitaire de bases. Le registre d’une occupation conserve le nombre de fois où intervient chaque irréductible.

Dans le secteur d'occupation totale un, \(\sum_p k_p=1\), de base \(\{e_p:p\in\mathcal P\}\), définissons

\[
H_{\mathcal P}e_p=(\log p)e_p,
\]
et, sur les occupations,

\[
H_{\mathrm{occ}}e_{\mathbf k}
=\left(\sum_pk_p\log p\right)e_{\mathbf k}.
\]
Avec les domaines diagonaux maximaux, tous deux sont autoadjoints. Si $Qe_n=ne_n$, l'identification des bases satisfait

\[
U H_{\mathrm{occ}}U^*=\log Q.
\]
L'égalité est une conséquence de la composition multiplicative :

\[
\sum_pk_p\log p=\log N(\mathbf k).
\]
L'opérateur n'est pas utilisé pour sélectionner rétrospectivement les nombres premiers. Il reçoit les irréductibles déjà construits et convertit leur composition en une quantité additive. C'est la place précise de la représentation des occupations dans la généalogie de l'article.

\subsection{Une même structure d'occupation et deux niveaux de trace}
Pour $\Re s>1$, la trace sur le secteur d'occupation totale un est

\[
P(s)=\operatorname{Tr}(e^{-sH_{\mathcal P}})=\sum_{p\in\mathcal P}p^{-s}.
\]
La trace sur toutes les occupations, en revanche, est

\[
Z(s)=\operatorname{Tr}(e^{-sH_{\mathrm{occ}}})
=\prod_{p\in\mathcal P}(1-p^{-s})^{-1}
=\sum_{n\ge1}n^{-s}.
\]
La dernière égalité provient de la factorisation unique. Le demi-plan indiqué garantit la convergence absolue et permet de regrouper les termes. Après l’identification de la réalisation entière, $Z$ y coïncide avec la fonction zêta. Aucune table de valeurs de zêta n’a été utilisée pour définir les modes.

\HMTresultado{Proposition}{Hiérarchie des moments}{rz-num-02-3}
Dans $\Re s>1$, en prenant la branche du logarithme déterminée par le développement du produit convergent,

\[
\log Z(s)=\sum_{k\ge1}\frac{P(ks)}{k}.
\]
\textbf{Démonstration.} Le développement absolument convergent $-\log(1-z)=\sum_{k\ge1}z^k/k$, appliqué à chaque $z=p^{-s}$, donne

\[
\log Z(s)=\sum_p\sum_{k\ge1}\frac{p^{-ks}}k.
\]
La convergence absolue autorise l'échange des deux sommes. \hfill\(\square\)

Les deux indices conservent des fonctions différentes : $p$ identifie le mode irréductible et $k$ sa multiplicité d'occupation. Leur relation ne se réduit pas à une liste de constantes : elle produit une famille de moments sur le même ensemble d'irréductibles.

\subsection{Meissel–Mertens et l'élimination de la contribution linéaire}
L'extension coordonnée réunit la partie finie

\[
B_1=\gamma+\sum_p\left(\log(1-p^{-1})+p^{-1}\right).
\]
Le terme ajouté $p^{-1}$ compense exactement le premier ordre du logarithme. Par conséquent,

\[
B_1=\gamma-\sum_{k\ge2}\frac{P(k)}k.
\]
Ici $\gamma$ désigne la partie finie $\gamma_0$ de la même fonctionnelle, dont la construction et la borne sont développées dans la section~\ref{viii:lectura:manuscrito-04} ; ce n’est ni une entrée métrologique ni une décimale choisie pour l’évaluation. La formule relie deux opérations sur la fonctionnelle construite : extraction de la partie finie et annulation du premier moment de chaque mode.

\HMTresultado{Proposition}{Borne de la queue des moments}{rz-num-02-4}
Pour un entier $K\ge1$,

\[
0<\sum_{k>K}\frac{P(k)}k
\le\frac{2P(K+1)}{K+1}
\le\frac{2^{-K}}{K+1}\left(1+\frac2K\right).
\]
\textbf{Démonstration.} Puisque $k\ge K+1$,

\[
\sum_{k>K}\frac{P(k)}k
\le\frac1{K+1}\sum_p\frac{p^{-(K+1)}}{1-p^{-1}}
\le\frac{2P(K+1)}{K+1}.
\]
Par comparaison avec tous les entiers supérieurs ou égaux à deux,

\[
P(K+1)\le\sum_{n\ge2}n^{-(K+1)}
\le2^{-(K+1)}+\int_2^\infty x^{-(K+1)}\,dx.
\]
L'évaluation de l'intégrale donne la seconde borne. \hfill\(\square\)

L'inégalité permet de convertir la somme finie en un intervalle, à condition que les évaluations finies et $\gamma$ possèdent à leur tour des bornes contrôlées. La longueur d'un développement décimal ne remplace pas cette estimation.

\subsection{Exclusion résiduelle et produit du motif de séparation deux}
Pour le motif à deux positions \(\{0,2\}\), soit $\nu_p$ le nombre de classes interdites modulo $p$. Dans $p=2$, les deux positions coïncident ; pour $p>2$, ce sont deux classes distinctes. La normalisation locale par rapport à deux choix indépendants produit

\[
\sigma_p=\frac{1-\nu_p/p}{(1-1/p)^2}.
\]
Le facteur $p=2$ vaut deux, tandis que les autres forment

\[
C_2=\prod_{p>2}\frac{1-2/p}{(1-1/p)^2}
=\prod_{p>2}\left(1-\frac1{(p-1)^2}\right).
\]
La série singulière complète du motif est $2C_2$. Il convient de réserver des symboles différents au produit sur $p>2$ et au facteur complet : la normalisation distingue une contribution résiduelle concrète, et non une convention facultative.

\HMTresultado{Proposition}{Le même système de moments avec un autre poids}{rz-num-02-5}
Le développement du produit satisfait

\[
\log C_2=-\sum_{k\ge2}\frac{2^k-2}{k}\bigl(P(k)-2^{-k}\bigr).
\]
\textbf{Démonstration.} Pour chaque $p>2$,

\[
\log(1-2/p)-2\log(1-1/p)
=-\sum_{k\ge2}\frac{2^k-2}{kp^k}.
\]
La somme commence à $k=2$, car le terme linéaire s’annule. Pour \(p\ge3\),

\[
-\log\left(1-\frac1{(p-1)^2}\right)\le\frac4{3(p-1)^2}.
\]
La comparaison avec la série des inverses des carrés démontre la finitude de la somme des valeurs absolues et permet d'en intervertir l'ordre. Le retrait du mode $2$ produit $P(k)-2^{-k}$. \hfill\(\square\)

La même hiérarchie $P(k)$ décrit ainsi deux opérations différentes : l’annulation linéaire de Meissel–Mertens et l’exclusion résiduelle du motif à deux positions. La différence réside dans le poids de chaque moment et dans le retrait du mode $2$, non dans l’introduction de deux catalogues indépendants de nombres premiers.

Pour un entier \(K\ge1\), soit

\[
S_K=\sum_{k=2}^K\frac{2^k-2}{k}\bigl(P(k)-2^{-k}\bigr),
\]
l'extension coordonnée donne la borne

\[
0<-\log C_2-S_K
\le\frac3{K+1}\left(\frac23\right)^{K+1}\left(1+\frac3K\right).
\]
Elle s'obtient en substituant $2^k-2\le2^k$, en sommant la série géométrique pour chaque $p\ge3$, en majorant $(1-2/p)^{-1}\le3$ et en comparant la somme sur les nombres premiers à la somme sur les entiers $n\ge3$. En prenant l'exponentielle, on obtient un intervalle pour $C_2$. La convergence de ce produit est un résultat distinct de l'existence d'une infinité de paires de nombres premiers séparés par deux.

\subsection{Réalisation triadique de la fonctionnelle}
Le propriétaire des fonctionnelles spectrales utilise les cycles $C_m=\mathbb Z/3^m\mathbb Z$ et leur laplacien

\[
(\Delta_mf)(j)=2f(j)-f(j+1)-f(j-1).
\]
En écrivant \(N_m=3^m\), soient \(e_{m,n}\) les caractères normalisés du cycle et \(\widetilde n\) le représentant de \(n\) entre \(-(N_m-1)/2\) et \((N_m-1)/2\). L'opérateur à orientation signée est défini par

\[
D_m e_{m,n}=\frac{N_m}{\sqrt{\pi}}
\sin\left(\frac{\pi\widetilde n}{N_m}\right)e_{m,n}.
\]
Ainsi est fixé le signe, et non seulement le carré. La réalisation satisfait

\[
D_m^2=\frac{3^{2m}}{4\pi}\Delta_m.
\]
Soit \(P_0\) le projecteur orthogonal sur les fonctions constantes du cycle. La coordonnée $\pi$ appartient à la génération antérieure du noyau HMT. Cette normalisation fait partie de la réalisation analytique et doit conserver sa provenance. La trace positive retire le mode constant et compte une contribution pour chaque couple d’orientations opposées :

\[
\operatorname{Tr}_m^+(F(D_m))
=\frac12\operatorname{Tr}\bigl(F(D_m)P_0^\perp\bigr).
\]
Avec les modes de Fourier du cycle, la trace de chaleur converge vers

\[
\Theta_+(x)=\sum_{n\ge1}e^{-\pi n^2x},\qquad x>0.
\]
La preuve, développée dans la première section du chapitre suivant, conserve une majorante gaussienne indépendante du niveau de raffinement. La transformée de Mellin de cette limite satisfait exactement

\[
\int_0^\infty\Theta_+(x)x^{s/2-1}\,dx
=\pi^{-s/2}\Gamma(s/2)Z(s),\qquad \Re s>1.
\]
Sont ainsi réunies la construction additive du spectre entier et la décomposition multiplicative par irréductibles. L'identité entre les fonctionnelles est démontrée ; elle ne se déduit pas de la proximité de deux évaluations numériques.

L'inversion thêta et le prolongement de la fonction complétée sont démontrés dans le chapitre suivant. Les parties finies, les coefficients de Laurent et les classes résiduelles modulo trois exigent leurs évaluations spécifiques ; ils ne se déduisent pas de la seule convergence de la trace de chaleur.

\subsubsection*{Défaut de transport et composante uniforme}
L’extension sur pli de feuillets fournit une représentation explicite du défaut d’un tour. Soit \(N=9^m\), avec \(m\ge1\), et soit \(T\) un opérateur unitaire sur une fibre de Hilbert \(\mathcal H\). Dans \(\mathbb C^N\), on utilise le produit scalaire euclidien usuel et le vecteur unitaire

\[
u_N=N^{-1/2}(1,\ldots,1).
\]
On définit \(\iota_Nf=u_N\otimes f\), \(P_N=\iota_N\iota_N^*\) et \(Q_N=I-P_N\). Le transport applique \(T\) une fois lorsque le cycle s'achève :

\[
U_N(T)=\sum_{r=1}^{N-1}|r+1\rangle\langle r|\otimes I
+|1\rangle\langle N|\otimes T.
\]
Sur l'état uniforme, il agit par

\[
U_N(T)\iota_N f
=\iota_N f+\frac{e_1}{\sqrt N}\otimes(T-I)f.
\]
Le terme ajouté possède une composante uniforme et une autre orthogonale. La seconde, avec la normalisation du propriétaire, est

\[
\delta_N(T)=\frac{N}{\sqrt{N-1}}Q_NU_N(T)\iota_N
=\eta_N\otimes(T-I),
\]
où

\[
\eta_N=\sqrt{\frac N{N-1}}\left(e_1-\frac{u_N}{\sqrt N}\right),
\qquad \|\eta_N\|=1,\qquad \eta_N\perp u_N.
\]
L'égalité se vérifie en substituant l'action précédente et en appliquant \(Q_N\). Pour deux transports unitaires \(T,S\) sur la même fibre, il vient

\[
\delta_N(T)^*\delta_N(S)=(I-T)^*(I-S).
\]
La commutativité n'est pas nécessaire : les produits apparaissent dans l'ordre indiqué. La même décomposition fournit

\[
\delta_N(T)^*P_N\delta_N(S)=0,
\]
puisque \(P_N\) annule le facteur \(\eta_N\).

Ces identités précisent l’information conservée par la normalisation du défaut et sa position par rapport à la composante uniforme. Elles permettent également de transporter isométriquement les fermetures des images entre deux longueurs de cycle : le gramien est indépendant de \(N\). Son domaine est le secteur du défaut construit ; il ne constitue pas à lui seul une identification de toutes les histoires du TPK ni de toute projection spectrale ultérieure.

\subsection{La forme complétée et son identité de frontière}
Le passage au double cercle utilise une forme sur un domaine de fonctions tests. Son côté arithmétique réunit les contributions premier–puissance, gamma et polaire. L'écriture à deux colonnes,

\[
\mathscr I_{\mathrm{GW}}=J_+^*J_+-J_-^*J_-,
\]
expose une différence de formes. Le producteur de l'indice premier, les poids orbitaux, la normalisation de Fourier et la complétion doivent figurer avant de l'identifier à une autre forme.

Le propriétaire du chapitre 96 situe l'étape centrale dans un codeur $\Xi$ et une projection orthogonale $P$, au moyen de

\[
\langle\Xi f,\Xi g\rangle=\langle J_+f,J_+g\rangle,
\qquad
\langle P\Xi f,P\Xi g\rangle=\langle J_-f,J_-g\rangle.
\]
Si ces deux identités s'obtiennent à partir des actions natives déclarées, leur différence donne exactement

\[
\mathscr I_{\mathrm{GW}}(f,g)
=\langle(I-P)\Xi f,(I-P)\Xi g\rangle.
\]
Le chapitre suivant démontre la décomposition explicite par colonnes et la conséquence des deux moments. L'application au codeur natif conserve comme antécédent son action sur les fonctions tests, la position de ses cartes par rapport à $P$ et la preuve des deux moments. L'orthogonalité de $P$ explique la dernière égalité, mais ne détermine pas à elle seule les deux identités précédentes.

Le télescopage du transport conserve pour sa part la forme

\[
\|x\|^2=\sum_{r=0}^{N-1}\|DA^rx\|^2+\|A^Nx\|^2,
\]
lorsque $D^*D=I-A^*A$. Si, sur le feuillet considéré, $A=qV$, avec $q=5/9$ et $V$ isométrique, le terme final vaut $q^{2N}\|x\|^2$ et tend vers zéro. La preuve de cette extinction est conservée avec l’identité qui situe l’état $x$ sur ce feuillet.

\subsection{Cercle spectral et cercle de mémoire}
Le premier cercle est une coordonnée de la droite spectrale. Pour $s=1/2+i\gamma$, avec \(\gamma\in\mathbb R\), la transformation

\[
\tau(s)=\frac{s-1}{s}
=\frac{\gamma+i/2}{\gamma-i/2}
\]
appartient au cercle unité privé du point \(1\) et s'inverse par

\[
s=\frac1{1-\tau},\qquad
\gamma=\frac12\cot\frac{\theta}{2},\quad\tau=e^{i\theta}.
\]
Le second cercle est la famille des caractères de la mémoire entière :

\[
\chi_\tau(n)=\tau^n,\qquad n\in\mathbb Z.
\]
La connexion nonadique augmente la mémoire d’une unité lorsqu’un tour s’achève. L’entrelacement relie une phase spectrale à son action sur cette mémoire ; il n’identifie pas le zéro numéro $j$ au tour numéro $j$.

Si \(C\in\mathcal B(\mathcal H)\) est l'opérateur de phase obtenu dans la réalisation spectrale, le calendrier à neuf positions agit sur \(\mathbb C^9\otimes\mathcal H\) par

\[
S_{C,9}=\sum_{r=1}^8|r+1\rangle\langle r|\otimes I
+|1\rangle\langle9|\otimes C.
\]
Chaque vecteur traverse une seule fois la liaison contenant $C$ au cours d'un tour complet. Par conséquent,

\[
S_{C,9}^9=I_9\otimes C.
\]
L'identité est exacte pour tout $C$ ; lorsque $C$ est unitaire, $S_{C,9}$ l'est également. Le facteur contractant $5/9$ appartient à une autre lecture du transport et se distingue de cette phase unitaire.

L'intérêt du chapitre ne réside pas dans la représentation d'une matrice quelconque sur neuf positions. Il réside dans la réunion de l'origine de l'opérateur $C$, de son domaine, du codeur précédent et de la compatibilité avec les cylindres du TPK. Ces données fixent ce qui est transporté et pourquoi cette réalisation est pertinente.

\clearpage
\subsection{Compatibilité avec la structure conjointe du continu}
L’espace des histoires, les deux feuillets, la mémoire, l’incidence et le complexe ne sont pas reconstruits comme des objets indépendants pour chaque section. L’assemblage corrélatif préalable est développé dans la section~\ref{viii:lectura:manuscrito-05}, avec ses opérations, ses conditions de réalisation et ses carrés de naturalité. Les opérateurs entier, de profondeur, d’irréductibles et de hauteurs spectrales restent typés ; une coïncidence de notation ne les identifie pas.

La section~\ref{viii:lectura:manuscrito-03} réunit l'identité de frontière avec sa construction terminale, les fonctions tests, les termes de la forme complétée et la convergence. Les approximations matricielles conservent leurs trois matrices de forme, de moment et de résidu. Les chiffres ne sont comparés qu'après la production et le gel des sorties.

Cette organisation conserve l'argument demandé : construction, opération, structure résultante et reconnaissance. Le manuscrit ne modifie pas le noyau causal pour faciliter un calcul ultérieur et n'utilise pas de liste de zéros comme entrée du producteur arithmétique.
